\section{Comparing results across model sizes and families}\label{sec:appendix-cross-model}

The main supervised results (\Cref{sec:supervised}) were produced on Llama-3.1-8B-Instruct. To check whether the persona LoRA recipe generalises beyond a single base model, we re-ran the pipeline (\Cref{sec:training-loras}) on three additional base models, holding the constitution, distillation procedure, and training hyper-parameters fixed: Gemma-3 4B/12B/27B-IT \citep{gemma_2025}, and Qwen2.5-7B-Instruct \citep{qwen2024qwen25}.

% TODO: verify caption matches final figure
We focus on the conscientiousness suppressor as the single trait/direction sweep, and evaluate each resulting adapter under the same Qwen3-235B judge (\Cref{sec:appendix-e}) on the conscientiousness rollout set, scaling the LoRA scalar over $\{-2, -1, 0, +1, +2\}$.

\Cref{fig:appendix-cross-model} shows the trait shift on the left and a coherence judge on the right, one line per base model. Two qualitative patterns are visible.
First, the LoRA recipe transfers across families and sizes: every base model moves in the intended (low-conscientiousness) direction at the trained scale and recovers approximately to baseline at scale $0$.
Second, the smaller Gemma checkpoints (4B and 12B) lose coherence noticeably as the suppressor is scaled up, whereas Gemma-3-27B and Llama-3.1-8B-Instruct remain near their baseline coherence over the same range. We read this as a capacity effect rather than a family effect: the smallest models trade general fluency for trait expression more readily, while larger checkpoints absorb the same intervention with less collateral damage.

% Figures removed: superseded by the cross-model OCEAN transfer matrices
% (appendices/model_comparison_matrices.tex) and the teacher ablation
% (appendices/teacher_ablation.tex). Prose above retained pending rewrite/merge;
% the \Cref{fig:appendix-cross-model} on the paragraph above is now dangling.

\FloatBarrier